\chapter*{Methodological Commitments}
\addcontentsline{toc}{chapter}{Methodological Commitments}
\label{ch:method_commitments}

\epigraph{A recurring mathematical structure licenses transfer of operational insight
only insofar as the transferred structure preserves measurable constraints, exclusions,
and failure conditions.}{}

\noindent This monograph 	tep{flyxion2025collapse} ranges across cosmology, cognition, computation, institutional
analysis, and paleoarchaeology. That range creates a specific risk: that recurring
mathematical structures will be taken as evidence of shared ontology rather than as
evidence that the same formal tools apply usefully across different substrates. This
chapter states the methodological commitments that guard against that risk.

\section*{The Analogy--Constraint Distinction}

\begin{firewall}
A recurring mathematical structure licenses transfer of operational insight only
insofar as the transferred structure preserves \emph{measurable constraints,
exclusions, and failure conditions}. Shared formalism is not evidence of identical
ontology. It is evidence that the same class of constraints operates at the relevant
scale.
\end{firewall}

This principle is the methodological spine of the monograph. Whenever a formal
framework developed in one domain is applied to another---when sheaf cohomology
appears in both semantic analysis and social institution modeling, or when the
\RSVP entropy field appears in both cosmological and cognitive contexts---the transfer
is licensed only by demonstrating that the relevant constraints, failure modes, and
exclusion conditions are structurally analogous. The transfer is not licensed by
aesthetic resonance, terminological similarity, or the intuition that the domains
``feel'' related.

The distinction is between \emph{analogy} and \emph{operational constraint transfer}.
An analogy asserts resemblance without specifying which properties are preserved and
which are not. An operational constraint transfer specifies exactly which constraints
are shared, which are not, and what would falsify the claim that the shared constraints
are sufficient to ground the transfer. Throughout this monograph, cross-domain
applications are intended as constraint transfers rather than analogies.

\section*{Claim-Status Architecture}

Because the monograph spans formally derived results, phenomenological interpretations,
engineering heuristics, conjectures, and speculative extrapolations, it uses a
consistent claim-status system to distinguish them. The five claim-status markers are:

\begin{claimstatus}
\textbf{[D] Definition.} A formal stipulation establishing terminology or a formal
object. Definitions are not true or false but more or less useful. They can be
evaluated by whether they generate productive inferences.

\textbf{[P] Phenomenological interpretation.} An application of a formal structure
to describe observable patterns without claiming that the formal structure is the
underlying mechanism. Phenomenological interpretations can be accurate descriptions
while remaining silent on mechanism.

\textbf{[H] Heuristic principle.} An operational guideline supported by evidence and
useful in practice, but not derived from first principles and potentially wrong in
edge cases. Heuristics should be tested against their domain of application.

\textbf{[C] Conjecture.} A claim that appears plausible given the available evidence
and the internal logic of the framework, but has not been formally derived or
empirically confirmed. Conjectures are productive when they generate testable
predictions.

\textbf{[S] Speculative extrapolation.} A claim that goes beyond what the available
evidence and formal derivations support, offered as a direction for investigation
rather than an established result.
\end{claimstatus}

These markers appear in the text at the points where the claim-status of a statement
shifts. They are not exhaustive annotations of every sentence: formally derived
theorems and definitions carry their status implicitly. The markers appear where
there is meaningful risk of status confusion---where a phenomenological
interpretation might be read as a formal claim, or where a conjecture might be
mistaken for an established result.

The most important risk is the silent transition from \textbf{[H]} to \textbf{[P]}
to \textbf{[D]}: a heuristic becomes a phenomenological description becomes a
formal definition without any single step being flagged as an escalation. The
claim-status system exists to make those transitions visible.

\section*{Constraint Boundaries}

Each major chapter ends with a \emph{Constraint Boundary} section that states
explicitly: what the framework in that chapter explains; what it does not explain;
what remains speculative; what would falsify the model; and which claims are
phenomenological versus formal. These sections are intended as internal audits
rather than summaries. They expose the limits of each local framework before the
next chapter extends it.

The monograph consistently argues against opaque abstraction---against systems
that compress away their dependency structure and present only terminal conclusions.
The Constraint Boundary sections are the mechanism by which the monograph resists
applying that pressure to itself.

\section*{On the RSVP Cosmological Material}

One specific clarification is required at the outset regarding the \RSVP cosmological
framework. The entropic relaxation account of cosmological redshift and the
admissibility geometry approach to structure formation are presented as
\textbf{[C]} conjectures with explicit falsification criteria, not as established
physics or replacements for the standard cosmological model. The redshift formula
$z_{\rm RSVP} \approx \int_0^L \kappa |\nabla S(\ell)| d\ell$ is a prediction of
the \RSVP field dynamics under specific coupling assumptions; it makes distinct
predictions at $k \gtrsim 0.1\,h/\text{Mpc}$ and at the Silk damping scale that
could in principle falsify the framework. Until those predictions are tested, the
cosmological material functions as a geometric field formalism for representing
structured evolution under entropy gradient, not as a replacement for established
observational cosmology.

\cleardoublepage
